\chapter{Loops, Goals, and Iterations: Agents, Skills, and Hooks}
\label{ch:loops-goals-iterations}

\begin{remark}[Learning Objectives]
After completing this chapter, the reader should be able to:
\begin{enumerate}
  \item Explain why one-shot prompting is fragile for complex work and describe
        the goal-directed loop (observe, propose, act, evaluate, revise) that
        replaces it.
  \item Write goals as concrete, checkable, decomposable control objects, and
        encode them in project files (TASK.md, TOPIC.md, STATUS.md, rubrics,
        validation scripts).
  \item Distinguish the roles of agents, skills, and hooks in an iterative
        system, and assemble them into a working iteration loop with explicit
        stopping criteria.
  \item Trace a deliberately simple feedback loop (derivative-free root finding)
        and a complex multi-agent loop (credit-risk stress testing) and explain
        why each requires real, feedback-driven iteration.
  \item Diagnose common failure modes of agentic iteration (infinite loops,
        fake convergence, goal drift, Goodhart overfitting, premature success)
        and apply hooks and stopping rules as remedies.
  \item Select an appropriate design pattern (planner--executor--critic,
        generate--test--revise, scaffold--draft--audit--repair, and others) for
        a given finance workflow.
\end{enumerate}
(This is a methodology capstone. The concrete tooling for Claude Code, Codex,
Hugging Face, the SDKs, Git, and Cline lives in Appendices~A--F; this chapter
treats the \emph{design of iterative agentic systems} rather than any single
tool.)
\end{remark}

% =====================================================================
\section{From Linear Prompts to Iterative Systems}
\label{sec:from-linear-to-iterative}
% =====================================================================

[Placeholder] Motivate the shift from single-shot prompting to goal-directed
iteration. Set up Claude Code (agents, skills, hooks, tests, validation scripts,
commits) as the concrete substrate, and ground the loop in the published
agentic-AI literature.

\subsection{Why One-Shot Prompting Is Fragile}
\label{subsec:one-shot-fragile}
[Placeholder] Explain how a single prompt has no error-correction channel:
errors compound, the model cannot see the consequences of its output, and
complex specifications exceed what one pass can satisfy. Contrast with how human
experts work iteratively.

\subsection{The Goal-Directed Loop: Observe, Propose, Act, Evaluate, Revise}
\label{subsec:goal-directed-loop}
[Placeholder] Define the five-phase loop. Each phase consumes the output of the
previous one; the loop only terminates when an explicit goal condition is met.

\subsection{Loops Versus Ordinary Repetition}
\label{subsec:loops-vs-repetition}
[Placeholder] Stress the distinction: repeating an action is not a loop unless
each iteration \emph{uses feedback} to change the next action. Give an example of
spurious repetition (re-prompting without reading the result) versus genuine
iteration.

\subsection{Claude Code as the Concrete Substrate}
\label{subsec:claude-code-substrate}
[Placeholder] Map the abstract loop onto concrete Claude Code mechanisms: agents,
skills, hooks, tests, validation scripts, and commits. Point readers to
Appendices~A--F for the tooling details; keep this section at the level of design.

\subsection{Situating the Loop in the Agentic-AI Literature}
\label{subsec:agentic-literature}
[Placeholder] Connect observe--propose--act--evaluate--revise to published
patterns: ReAct (reason + act), Reflexion (self-feedback), Self-Refine, and the
planner--executor--critic decomposition. Establish that the method is grounded,
not bespoke to one CLI.

\subsection{When Not to Use an Agentic Loop}
\label{subsec:when-not-to-loop}
[Placeholder] Pre-empt over-engineering: a plain script, a pure function, or a
cron job often beats an agent. Give a decision heuristic for when iteration earns
its cost. Pairs with the ``excessive agent delegation'' failure mode in
Section~\ref{sec:failure-modes}.

% =====================================================================
\section{Goals as Control Objects}
\label{sec:goals-as-control}
% =====================================================================

[Placeholder] Treat the goal as the object that \emph{controls} the loop: it
determines what counts as progress and when to stop.

\subsection{Concrete, Checkable, Decomposable Goals}
\label{subsec:concrete-checkable-goals}
[Placeholder] A usable goal must be (i) concrete (states the artifact), (ii)
checkable (a machine or rubric can verify it), and (iii) decomposable (breaks
into local subgoals). Show how each property maps to a loop mechanism.

\subsection{Success Criteria, Stopping Rules, and Failure Modes}
\label{subsec:success-stopping-failure}
[Placeholder] Distinguish the success criterion (what ``done'' means) from the
stopping rule (when to halt regardless of success) and the anticipated failure
modes (what to watch for). Every goal needs all three.

\subsection{Representing Goals in Files}
\label{subsec:goals-in-files}
[Placeholder] Show how goals live in TASK.md, TOPIC.md, STATUS.md, rubric files,
and validation scripts. A goal in a file is durable across context windows and
auditable; a goal only in the prompt is lost. Reference this book's own
\texttt{TOPIC.md} and \texttt{docs/quality/RUBRIC.md}.

\subsection{Bad Goals Versus Good Goals}
\label{subsec:bad-vs-good-goals}
[Placeholder] Worked contrasts. Bad: ``improve the chapter.'' Good: ``raise every
rubric dimension to $\geq 90$ as measured by \texttt{/score-content}.'' Show a
table of three or four bad/good pairs with the reason each fails or works.

% =====================================================================
\section{Agents as Specialized Iterators}
\label{sec:agents-as-iterators}
% =====================================================================

[Placeholder] Frame agents as specialized workers inside the loop, not as
general chatbots. Keep at the level of design; how to author an agent file is
appendix material.

\subsection{Roles: Worker, Critic, Validator, Planner, Refactorer}
\label{subsec:agent-roles}
[Placeholder] Catalogue the standard roles an agent can play and what each
contributes to an iteration.

\subsection{Single-Agent Versus Multi-Agent Loops}
\label{subsec:single-vs-multi-agent}
[Placeholder] When one agent suffices and when separation of concerns (e.g., a
generator and an independent critic) materially improves outcomes.

\subsection{Handoffs and Shared State Between Agents}
\label{subsec:agent-handoffs}
[Placeholder] How agents pass work to one another: shared files, structured
outputs, and commits as the medium of handoff. The danger of lossy handoffs.

\subsection{Agents Should Improve Project State, Not Merely Emit Output}
\label{subsec:agents-improve-state}
[Placeholder] The key discipline: an agent's value is the change it makes to the
durable project state (files, tests, scores), not the text it prints.

% =====================================================================
\section{Skills as Reusable Procedures}
\label{sec:skills-as-procedures}
% =====================================================================

[Placeholder] Frame skills as the reusable, composable procedures of the system.

\subsection{Skills as Reusable Operational Knowledge}
\label{subsec:skills-operational-knowledge}
[Placeholder] A skill captures ``how we do X here'' once, so it need not be
re-derived each session.

\subsection{Skills Compared to Functions}
\label{subsec:skills-as-functions}
[Placeholder] The load-bearing analogy: skills are like functions --- scoped,
documented, repeatable, and composable. Map each function property to a skill
property and show a skill calling another skill.

\subsection{Skills Encode Project-Specific Conventions}
\label{subsec:skills-encode-conventions}
[Placeholder] How a skill freezes a project's conventions (numbering, commit
format, quality gates) so every run is consistent. Use this book's
\texttt{/new-topic} and \texttt{/draft-chapter} as examples.

\subsection{Skills Reduce Entropy in Long-Running Projects}
\label{subsec:skills-reduce-entropy}
[Placeholder] Over a long project, ad hoc actions accumulate inconsistency;
skills are the counterforce that keeps the project ordered.

% =====================================================================
\section{Hooks as Guardrails and Feedback Channels}
\label{sec:hooks-as-guardrails}
% =====================================================================

[Placeholder] Frame hooks as automatic interventions that convert subjective
``is this good?'' into measurable, enforced feedback.

\subsection{Hooks as Automatic Interventions}
\label{subsec:hooks-automatic}
[Placeholder] Define pre-action and post-action hooks and how they fire without
the agent's cooperation.

\subsection{The Hook Catalogue}
\label{subsec:hook-catalogue}
[Placeholder] Pre-run checks; post-write validation; linting, tests,
compilation; safety checks; and commit checks. One concrete finance example per
type.

\subsection{Turning Subjective Improvement into Measurable Feedback}
\label{subsec:hooks-measurable}
[Placeholder] How a hook gives the loop a crisp pass/fail signal, replacing the
model's self-assessment with an external one.

\subsection{Designing Trustworthy Validators}
\label{subsec:trustworthy-validators}
[Placeholder] A passing check is not the same as a met goal. Discuss how to
build validators that are hard to game, setting up the Goodhart failure mode in
Section~\ref{sec:failure-modes}.

\subsection{Failure Hooks}
\label{subsec:failure-hooks}
[Placeholder] What should happen when a test fails, a build breaks, or a rubric
score drops: block the commit, revert, escalate, or re-enter the loop.

% =====================================================================
\section{Designing an Iteration Loop}
\label{sec:designing-loop}
% =====================================================================

[Placeholder] Synthesize Sections~\ref{sec:goals-as-control}--\ref{sec:hooks-as-guardrails}
into a concrete, runnable loop design.

\subsection{The General Loop Template}
\label{subsec:loop-template}
[Placeholder] The seven steps: read current state; state the next local goal;
propose a change; apply the change; run validation; interpret feedback; decide to
stop, continue, or escalate. Present as a numbered procedure.

\subsection{Pseudocode and Its Mapping onto a Claude Code Workflow}
\label{subsec:loop-pseudocode}
[Placeholder] Give clean pseudocode for the loop, then annotate how each line
maps onto agents, skills, hooks, and commits in practice.

\subsection{State, Memory, and Context Management Across Iterations}
\label{subsec:state-memory-context}
[Placeholder] What persists between turns (files, commits, status tables) versus
what is lost to the context window. Practical strategies for carrying state
forward without relying on conversation memory.

\subsection{Stopping Criteria}
\label{subsec:stopping-criteria}
[Placeholder] Convergence (no further improvement), budget exhaustion,
validation pass, and human-review-required. Every loop needs at least one
guaranteed-terminating criterion.

\subsection{The Economics of Iteration}
\label{subsec:economics-of-iteration}
[Placeholder] Cost per iteration, latency, and token/compute budgets as
first-class stopping criteria. The ROI question: is this loop worth running at
all? Particularly relevant for production finance workflows.

% =====================================================================
\section{Example 1: A Naive Derivative-Free Root-Finding Loop}
\label{sec:example-root-finding}
% =====================================================================

[Placeholder] A deliberately simple worked example. The pedagogical point is not
numerical methods but that an agentic loop can improve a state variable using
feedback alone.

\subsection{Problem Setup: A Zero of $f(x)=x^3-x-2$ Without Derivatives}
\label{subsec:rootfind-setup}
[Placeholder] State the problem: find $x$ with $f(x)=x^3-x-2 \approx 0$. The
agent may evaluate $f$ but is forbidden from using derivatives, so this is
Newton-\emph{inspired} only in the loose sense of iterative approximation. Note
the finance analogue: the same mechanics solve for an implied yield or internal
rate of return by sign-guided search.

\subsection{The Supporting Artifacts: Goal File, Skill, and Hook}
\label{subsec:rootfind-artifacts}
[Placeholder] A goal file defining the tolerance $|f(x)| < 0.001$; a skill
describing how to run the evaluator script; and a hook that prevents the agent
from declaring success unless the tolerance is satisfied.

\subsection{A Realistic Iteration Trace from a Poor Start}
\label{subsec:rootfind-trace}
[Placeholder] Start from $x=0$ so several real iterations are needed. Show the
trace where the agent reads the sign and magnitude of the error and proposes a
better guess, e.g.:
$x=0,\ f=-2$; $x=2,\ f=4$; $x=1.5,\ f=-0.125$; $x=1.6,\ f=0.496$;
$x=1.52,\ f=-0.008$; $x=1.522,\ f\approx 0.004$; $x=1.5214,\ f\approx 0$.

\subsection{The Final Report}
\label{subsec:rootfind-report}
[Placeholder] The loop emits a short report summarizing the iterations, the final
$x$, the final $|f(x)|$, and confirmation that the tolerance hook passed.

% =====================================================================
\section{Example 2: A Credit-Risk Stress-Testing Agent Loop}
\label{sec:example-credit-risk}
% =====================================================================

[Placeholder] The capstone case study. A mid-sized private company applies for a
large revolving credit facility; it looks healthy in the base case but may
deteriorate under a combined macro-financial shock. The object being optimized is
not a single number but a defensible credit judgment that must be internally
consistent, empirically grounded, stress-tested, and survive skeptical review ---
which is exactly why it requires many feedback-driven iterations.

\subsection{The Goal: A Defensible Credit-Committee Memo}
\label{subsec:credit-goal}
[Placeholder] Central goal: produce a memo assigning a probability of default
(PD), loss-given-default (LGD) assumptions, expected loss, key risk drivers, and
a final lend / do-not-lend recommendation that a credit committee would accept.

\subsection{The Combined Macro-Financial Shock}
\label{subsec:credit-shock}
[Placeholder] Define the combined stress: revenue decline, margin compression,
higher interest rates, delayed customer payments, inventory buildup, covenant
tightening, and refinancing risk. Each must enter the model with explicit
assumptions.

\subsection{The Agent Ensemble}
\label{subsec:credit-agents}
[Placeholder] Specialized agents and their handoffs: data-ingestion,
accounting-quality, feature-engineering, PD-model, stress-scenario, cash-flow,
LGD, skeptic, and memo-writing. Describe what each contributes and what it hands
to the next.

\subsection{The Skills}
\label{subsec:credit-skills}
[Placeholder] Reusable procedures invoked by the agents: financial-statement
normalization; credit-feature engineering; probability-of-default modeling; macro
stress-scenario design; covenant and liquidity analysis; LGD and collateral
valuation; and credit-committee memo writing.

\subsection{The Hooks}
\label{subsec:credit-hooks}
[Placeholder] Guardrails enforced throughout: no uncited factual claims; all
ratios reconcile to source statements; train/test separation before evaluation;
target-leakage checks pass before any PD is reported; every scenario specifies
assumptions in a parameter table; the recommendation is not final until the
skeptic agent signs off; code regenerates all tables, figures, and exhibits
end-to-end; and the memo may not state ``low risk'' unless both base-case and
stressed expected loss are below explicit thresholds.

\subsection{The Eleven-Iteration Narrative}
\label{subsec:credit-iterations}
[Placeholder] The loop in motion:
(1) naive memo from historical profitability and leverage;
(2) accounting-quality agent finds EBITDA add-backs and inconsistent
working-capital definitions;
(3) feature-engineering agent adds payment delays, customer concentration, cash
conversion cycle, and covenant headroom;
(4) PD-model agent's first model is rejected when leakage hooks fire;
(5) model rebuilt on clean features with out-of-sample validation;
(6) stress-scenario agent imposes the combined shock (revenue $-20\%$, gross
margin $-300$\,bps, rates $+250$\,bps, DSO $+25$ days, slower inventory turnover);
(7) cash-flow agent finds the borrower survives the base case but breaches
interest-coverage and minimum-liquidity covenants under stress;
(8) LGD agent finds weaker recovery because inventory is specialized and
receivables are concentrated in two customers;
(9) skeptic agent demands alternative macro assumptions, peer benchmarking, and
recovery-rate sensitivity;
(10) agents rerun PD, LGD, expected loss, and covenants under base, adverse, and
severe scenarios;
(11) memo-writing agent produces the final recommendation (reject, smaller
facility, tighter covenants, additional collateral, or pricing compensation).

\subsection{The Compact Final Result Table}
\label{subsec:credit-table}
[Placeholder] A compact summary table with base-case PD, stressed PD, base-case
LGD, stressed LGD, expected loss, covenant-breach indicator, liquidity shortfall,
and recommended action.

% =====================================================================
\section{Human Oversight, Accountability, and Audit Trails}
\label{sec:human-oversight}
% =====================================================================

[Placeholder] Distinct from hooks: in regulated finance the loop's output must be
signed off, attributable, and reproducible. The credit-memo example raises
``who is accountable?''; this section answers it.

\subsection{Human-in-the-Loop and Escalation Points}
\label{subsec:human-in-the-loop}
[Placeholder] Where a human must approve before the loop proceeds, and how
escalation is triggered.

\subsection{Accountability and Sign-Off in Regulated Finance}
\label{subsec:accountability-signoff}
[Placeholder] Model-risk governance: who owns the model, who validates it, and
how sign-off is recorded.

\subsection{Reproducible Audit Trails}
\label{subsec:audit-trails}
[Placeholder] Commits, logs, and iteration traces as evidence. An audit-ready
loop can reconstruct exactly how every number in the final memo was produced.

% =====================================================================
\section{Failure Modes in Agentic Iteration}
\label{sec:failure-modes}
% =====================================================================

[Placeholder] Catalogue the ways agentic loops fail, with warning signs and
remedies. Emphasize that hooks and explicit stopping rules are the primary
defenses.

\subsection{Infinite Loops and Fake Convergence}
\label{subsec:infinite-fake-convergence}
[Placeholder] Loops that never stop, and loops that declare convergence without
real progress. Warning signs and remedies.

\subsection{Goal Drift and Premature Success}
\label{subsec:goal-drift-premature}
[Placeholder] The goal silently mutating across iterations, and the agent
declaring success before the goal is met. How file-based goals and tolerance
hooks prevent both.

\subsection{Overfitting to Validation Checks (Goodhart's Law)}
\label{subsec:goodhart-overfitting}
[Placeholder] When a measure becomes a target it ceases to be a good measure: the
agent games the validator instead of meeting the goal. Remedies: diverse checks,
held-out validation, adversarial review.

\subsection{Excessive Agent Delegation}
\label{subsec:excessive-delegation}
[Placeholder] Spawning agents where a function would do; coordination overhead
swamping benefit. Connects back to Section~\ref{subsec:when-not-to-loop}.

\subsection{Warning Signs and Remedies}
\label{subsec:warning-signs-remedies}
[Placeholder] A consolidated table mapping each failure mode to its early warning
signs and the hook or stopping rule that mitigates it.

% =====================================================================
\section{Practical Design Patterns}
\label{sec:design-patterns}
% =====================================================================

[Placeholder] A catalogue of reusable loop architectures and when each applies.

\subsection{Planner--Executor--Critic}
\label{subsec:pattern-planner-executor-critic}
[Placeholder] One agent plans, one executes, one critiques. When separation of
planning from execution pays off.

\subsection{Generate--Test--Revise}
\label{subsec:pattern-generate-test-revise}
[Placeholder] The tightest feedback loop: produce, validate, fix, repeat. The
default for well-specified, machine-checkable goals.

\subsection{Scaffold--Draft--Audit--Repair}
\label{subsec:pattern-scaffold-draft-audit-repair}
[Placeholder] The pattern behind this very book: scaffold structure, draft
content, audit against a rubric, repair failures. When goals are subjective and
rubric-scored.

\subsection{Data-Pipeline: Build--Validate--Document}
\label{subsec:pattern-data-pipeline}
[Placeholder] For reproducible data work: build the pipeline, validate the
outputs, document the lineage. When provenance matters.

\subsection{Credit-Model: Build--Stress-Test--Red-Team--Finalize}
\label{subsec:pattern-credit-model}
[Placeholder] Generalize Example~2 (Section~\ref{sec:example-credit-risk}) into a
reusable pattern for any high-stakes financial model.

\subsection{Choosing a Pattern}
\label{subsec:choosing-a-pattern}
[Placeholder] A decision guide mapping properties of the goal (checkable vs.
subjective, single-number vs. judgment, stakes, reproducibility needs) to the
appropriate pattern.

% =====================================================================
\section{Summary and Takeaways}
\label{sec:loops-summary}
% =====================================================================

[Placeholder] Close the chapter.

\subsection{Iteration Is Useful Only When Feedback Is Informative}
\label{subsec:takeaway-feedback}
[Placeholder] The central thesis restated: a loop without an informative feedback
signal is just expensive repetition.

\subsection{Agents Specialize the Loop; Skills Reuse It; Hooks Enforce It}
\label{subsec:takeaway-three-blocks}
[Placeholder] Recap the division of labor: agents specialize parts of the loop,
skills make procedures reusable, and hooks make quality control automatic.

\subsection{Complex Finance Workflows Demand Explicit Goals, Validation, Stress Testing, and Skepticism}
\label{subsec:takeaway-finance}
[Placeholder] Why finance in particular needs this discipline, drawing the thread
from the credit-risk example through governance and audit trails.
